% ============================================================
% Laporan Akhir Modul 1 - Pengantar Mikrokontroler
% Penanda MERAH = harus diisi/dicek dulu sebelum dikumpulkan:
%   \isi{...}  -> teks yang perlu kamu lengkapi atau konfirmasi
%   [FOTO: nama_file] -> otomatis berubah jadi gambar bila file
%                        P1/img/nama_file.png (atau .jpg) ada
% ============================================================
\newcommand{\isi}[1]{\textcolor{red}{\textbf{[ISI: #1]}}}
% Bagian B Percobaan 3 (pull-up/pull-down INTERNAL). Ubah ke \BagianBtrue HANYA jika
% Bagian B benar-benar dijalankan dan ada screenshot keluaran [INTERNAL].
\newif\ifBagianB
\BagianBtrue
% Kotak merah untuk foto yang belum ada. Pada build FINAL ubah ke \ShowMissingfalse
% supaya figur yang filenya belum ada tidak ditampilkan sama sekali.
\newif\ifShowMissing
\ShowMissingfalse
\newcommand{\foto}[3][0.75]{%
  \IfFileExists{P1/img/#2.png}{\fotoada[#1]{P1/img/#2.png}{#3}}{%
    \IfFileExists{P1/img/#2.jpg}{\fotoada[#1]{P1/img/#2.jpg}{#3}}{%
      \ifShowMissing
        \begin{figure}[H]\centering
          \fbox{\parbox[c][4cm][c]{#1\linewidth}{\centering\textcolor{red}{\textbf{[FOTO: \detokenize{#2}]}}}}
          \caption{#3}
        \end{figure}
      \fi}}}
\newcommand{\fotoada}[3][0.75]{%
  \begin{figure}[H]
    \centering
    \includegraphics[width=#1\linewidth]{#2}
    \caption{#3}
  \end{figure}}

%===========================================================%
\section{Langkah-Langkah Percobaan}
Praktikum Modul 1 dilaksanakan pada 8 November 2026. Seluruh percobaan memakai satu project PlatformIO bernama \texttt{modul1-praktikum} di Visual Studio Code. Folder project diletakkan pada path tanpa spasi. Setiap percobaan hanya mengganti isi \texttt{platformio.ini} dan file kode di \texttt{src/}. Alat dan bahan mengikuti Modul 1: STM32 Blackpill, ST-Link, ESP32 DevKit, pushbutton, resistor 10~k$\Omega$, serta kabel jumper.

\subsection{Percobaan 1: STM32 Blackpill (Blink)}
\begin{enumerate}
	\item Mengisi \texttt{platformio.ini} dengan konfigurasi berikut, lalu menunggu PlatformIO selesai me-refresh konfigurasi.
\begin{lstlisting}[language=C, caption={platformio.ini Percobaan 1}]
[env:blackpill_f401cc]
platform = ststm32
board = blackpill_f401cc
framework = arduino
upload_protocol = stlink
debug_tool = stlink
\end{lstlisting}

	\item Menghubungkan ST-Link ke Blackpill sesuai Tabel~\ref{tab:stlink}, lalu menghubungkan ST-Link ke komputer melalui USB.
\begin{table}[H]
	\centering
	\caption{Sambungan ST-Link ke STM32 Blackpill}
	\label{tab:stlink}
	\begin{tabular}{|l|l|l|}
		\hline
		\textbf{Pin ST-Link} & \textbf{Pin Blackpill} & \textbf{Fungsi} \\ \hline
		SWDIO & SWDIO & Jalur data debug \\ \hline
		SWCLK & SWCLK & Jalur clock debug \\ \hline
		GND & GND & Ground bersama \\ \hline
		3.3V & 3V3 & Tegangan target \\ \hline
	\end{tabular}
\end{table}
\foto{p1_wiring_stlink}{Sambungan ST-Link ke STM32 Blackpill}

	\item Menulis program Blink pada \texttt{src/main.cpp}, lalu menjalankan Build dan Upload.
\begin{lstlisting}[language=C++, caption={Program Blink STM32 (Arduino)}]
#include <Arduino.h>

#define LED_PIN PC13

void setup()
{
  pinMode(LED_PIN, OUTPUT);
}

void loop()
{
  digitalWrite(LED_PIN, LOW);  // LED onboard aktif LOW
  delay(500);
  digitalWrite(LED_PIN, HIGH);
  delay(500);
}
\end{lstlisting}
\foto{p1_build_upload}{Hasil Build dan Upload Percobaan 1 pada terminal PlatformIO}

	\item Mengamati LED onboard pada pin PC13.
\foto{p1_led_pc13}{LED onboard PC13 pada STM32 Blackpill saat program berjalan}
\end{enumerate}
Upload pertama sempat gagal. Penyebab dan penanganannya dibahas pada bagian Analisis Hasil Percobaan.

\subsection{Percobaan 2: ESP32 dengan Framework ESP-IDF (Blink)}
\begin{enumerate}
	\item Mengganti isi \texttt{platformio.ini} dengan konfigurasi ESP-IDF.
\begin{lstlisting}[language=C, caption={platformio.ini Percobaan 2}]
[env:esp32dev]
platform = espressif32
board = esp32dev
framework = espidf
\end{lstlisting}

	\item Menghubungkan ESP32 ke komputer melalui USB dan memastikan port serialnya terdeteksi (pada kami COM5).
\foto[0.55]{p2_port_com5}{Device Manager: ESP32 terdeteksi sebagai USB-SERIAL CH340 (COM5)}
	\item Menghapus \texttt{src/main.cpp}, membuat \texttt{src/main.c}, menulis program berikut, lalu menjalankan Build dan Upload.
\begin{lstlisting}[language=C, caption={Program Blink ESP32 (ESP-IDF)}]
#include "freertos/FreeRTOS.h"
#include "freertos/task.h"
#include "driver/gpio.h"

#define LED_PIN GPIO_NUM_2

void app_main(void)
{
  gpio_reset_pin(LED_PIN);
  gpio_set_direction(LED_PIN, GPIO_MODE_OUTPUT);

  while (1)
  {
    gpio_set_level(LED_PIN, 1);
    vTaskDelay(pdMS_TO_TICKS(500));
    gpio_set_level(LED_PIN, 0);
    vTaskDelay(pdMS_TO_TICKS(500));
  }
}
\end{lstlisting}
\foto{p2_build_upload}{Hasil Build dan Upload Percobaan 2 pada terminal PlatformIO}

	\item Mengamati LED onboard ESP32 pada GPIO 2.
\foto{p2_led_gpio2}{LED onboard ESP32 (GPIO 2) berkedip}
\end{enumerate}

\subsection{Percobaan 3: Pull-up dan Pull-down, Eksternal dan Internal}
\textbf{Bagian A: resistor eksternal.}
\begin{enumerate}
	\item Mengembalikan \texttt{platformio.ini} ke framework Arduino, menghapus \texttt{src/main.c}, dan membuat kembali \texttt{src/main.cpp}. Pada percobaan ini kami menambahkan \texttt{monitor\_speed = 115200} agar sama dengan \texttt{Serial.begin(115200)} pada kode.
\begin{lstlisting}[language=C, caption={platformio.ini Percobaan 3}]
[env:esp32dev]
platform = espressif32
board = esp32dev
framework = arduino
monitor_speed = 115200
\end{lstlisting}

	\item Merangkai dua pushbutton: tombol pull-up pada GPIO 32 (3V3 melalui resistor 10~k$\Omega$ ke node GPIO 32, tombol dari node ke GND) dan tombol pull-down pada GPIO 33 (GND melalui resistor 10~k$\Omega$ ke node GPIO 33, tombol dari node ke 3V3).
\foto{p3_wiring_eksternal}{Rangkaian pull-up (GPIO 32) dan pull-down (GPIO 33) eksternal}

	\item Menulis Kode Program A, menjalankan Build dan Upload, membuka serial monitor, lalu menekan kedua tombol bergantian.
\begin{lstlisting}[language=C++, caption={Kode Program A: pull-up dan pull-down eksternal}]
#include <Arduino.h>

#define BUTTON_PULLUP_PIN 32   // pull-up eksternal
#define BUTTON_PULLDOWN_PIN 33 // pull-down eksternal

void setup() {
  Serial.begin(115200);
  pinMode(BUTTON_PULLUP_PIN, INPUT);
  pinMode(BUTTON_PULLDOWN_PIN, INPUT);
}

void loop() {
  int stateUp = digitalRead(BUTTON_PULLUP_PIN);
  int stateDown = digitalRead(BUTTON_PULLDOWN_PIN);

  Serial.printf("[EKSTERNAL] Pull-up (GPIO32): %s | Pull-down (GPIO33): %s\n",
                stateUp == LOW ? "DITEKAN" : "TIDAK DITEKAN",
                stateDown == HIGH ? "DITEKAN" : "TIDAK DITEKAN");
  delay(200);
}
\end{lstlisting}
\foto{p3_serial_eksternal}{Keluaran serial monitor Kode Program A setelah wiring diperbaiki (tombol tidak ditekan)}
\end{enumerate}

\ifBagianB
\textbf{Bagian B: resistor internal.}
\begin{enumerate}
	\item Melepas kedua resistor eksternal, tombol tetap terpasang pada GPIO 32 (kaki lain ke GND) dan GPIO 33 (kaki lain ke 3V3).
\foto{p3_wiring_internal}{Rangkaian tombol tanpa resistor eksternal (pull-up dan pull-down internal)}

	\item Mengganti isi \texttt{src/main.cpp} dengan Kode Program B, menjalankan Build dan Upload, lalu menekan kedua tombol.
\begin{lstlisting}[language=C++, caption={Kode Program B: pull-up dan pull-down internal}]
#include <Arduino.h>

#define BUTTON_PULLUP_PIN 32   // pull-up internal
#define BUTTON_PULLDOWN_PIN 33 // pull-down internal

void setup() {
  Serial.begin(115200);
  pinMode(BUTTON_PULLUP_PIN, INPUT_PULLUP);
  pinMode(BUTTON_PULLDOWN_PIN, INPUT_PULLDOWN);
}

void loop() {
  int stateUp = digitalRead(BUTTON_PULLUP_PIN);
  int stateDown = digitalRead(BUTTON_PULLDOWN_PIN);

  Serial.printf("[INTERNAL] Pull-up (GPIO32): %s | Pull-down (GPIO33): %s\n",
                stateUp == LOW ? "DITEKAN" : "TIDAK DITEKAN",
                stateDown == HIGH ? "DITEKAN" : "TIDAK DITEKAN");
  delay(200);
}
\end{lstlisting}
\foto{p3_serial_internal}{Keluaran serial monitor Kode Program B: tombol tidak ditekan, kemudian tombol pull-down (GPIO 33) ditekan}
\end{enumerate}
\else
Bagian B (pull-up dan pull-down internal) tidak sempat dilaksanakan, sehingga langkah dan hasilnya tidak dimuat dalam laporan ini.
\fi

\subsection{Percobaan 4 dan 5}
Percobaan 4 (debouncing) dan Percobaan 5 (ADC dengan LDR) tidak sempat dilaksanakan karena waktu praktikum habis setelah Percobaan 3. Oleh karena itu laporan ini tidak memuat langkah, data, maupun analisis hasil untuk kedua percobaan tersebut.

%===========================================================%
\section{Analisis Hasil Percobaan}

\subsection{Percobaan 1}
LED onboard PC13 menyala 500~ms lalu mati 500~ms secara berulang, sehingga periodenya 1000~ms (frekuensi 1~Hz) dengan duty cycle 50\%. Hasil ini sesuai dengan program: dua pemanggilan \texttt{delay(500)} dalam satu putaran \texttt{loop()}. LED onboard Blackpill bersifat aktif LOW, yaitu anodanya terhubung ke 3,3~V dan katodanya ke pin PC13, sehingga LED menyala ketika pin diberi LOW. Karena itu \texttt{digitalWrite(LED\_PIN, LOW)} menyalakan LED dan \texttt{HIGH} mematikannya, sesuai dengan teori pada modul.

Upload pertama gagal dengan pesan OpenOCD \texttt{Error: open failed}, dan Device Manager menampilkan \emph{Unknown USB Device (Device Descriptor Request Failed)}. Build sudah berhasil, jadi kesalahan bukan pada kode, melainkan pada komunikasi komputer dengan ST-Link melalui USB. Kami memasang driver ST-Link (paket \texttt{stsw-link009}), tetapi perangkat masih bertanda seru di Device Manager. Setelah sambungan kabel SWD diperiksa dan kesalahan wiring diperbaiki, upload berhasil. Pelajaran dari kendala ini adalah urutan pemeriksaan: kode (build berhasil atau tidak), driver, lalu sambungan fisik.
\foto[0.55]{p1_error_usb}{Device Manager menampilkan \emph{Unknown USB Device (Device Descriptor Request Failed)} saat upload pertama gagal}

\subsection{Percobaan 2}
LED onboard pada GPIO 2 berkedip dengan periode 1000~ms, sama seperti Percobaan 1, sehingga program ESP-IDF terbukti berjalan di ESP32. Karena sama-sama Blink, perbedaan Arduino dan ESP-IDF terlihat pada struktur programnya (Tabel~\ref{tab:arduino-idf}).
\begin{table}[H]
	\centering
	\caption{Perbandingan program Blink dengan Arduino dan ESP-IDF}
	\label{tab:arduino-idf}
	\begin{tabular}{|p{3.2cm}|p{5.6cm}|p{5.6cm}|}
		\hline
		\textbf{Aspek} & \textbf{Arduino} & \textbf{ESP-IDF} \\ \hline
		Titik masuk & \texttt{setup()} sekali, \texttt{loop()} berulang & \texttt{app\_main()}, perulangan ditulis sendiri \\ \hline
		Konfigurasi pin & \texttt{pinMode()} & \texttt{gpio\_reset\_pin()} dan \texttt{gpio\_set\_direction()} \\ \hline
		Menulis level pin & \texttt{digitalWrite()} & \texttt{gpio\_set\_level()} \\ \hline
		Penundaan & \texttt{delay()} (ms) & \texttt{vTaskDelay(pdMS\_TO\_TICKS())} \\ \hline
		Berkas utama & \texttt{main.cpp} & \texttt{main.c} \\ \hline
	\end{tabular}
\end{table}
ESP-IDF menuntut lebih banyak baris dan pemahaman API tingkat rendah, tetapi memberi akses langsung ke FreeRTOS dan fitur ESP32. Arduino lebih ringkas karena detail tersebut disembunyikan oleh framework. Pada ESP32, framework Arduino sendiri berjalan di atas FreeRTOS, sehingga perbedaan keduanya terutama pada tingkat abstraksi dan kontrol, bukan pada ada atau tidaknya RTOS.

Ada dua hal teknis yang memengaruhi keberhasilan build ESP-IDF di PlatformIO. Pertama, saat persiapan sebelum praktikum, project yang berada pada path dengan spasi gagal dikonfigurasi (\emph{Detected a whitespace character in project paths}), sehingga project dipindahkan ke path tanpa spasi. Kedua, program ESP-IDF harus disimpan sebagai \texttt{main.c}. Pada file C++ nama fungsi mengalami \emph{name mangling}, sehingga \texttt{app\_main} tidak dikenali kecuali dideklarasikan \texttt{extern~\textquotedbl C\textquotedbl}.

\subsection{Percobaan 3}
Tabel~\ref{tab:pullup} merangkum perilaku yang diharapkan menurut teori, dan keluaran serial monitor pada kode memetakannya menjadi tulisan DITEKAN atau TIDAK DITEKAN.
\begin{table}[H]
	\centering
	\caption{Level pin yang diharapkan pada Percobaan 3}
	\label{tab:pullup}
	\begin{tabular}{|l|l|l|}
		\hline
		\textbf{Kondisi tombol} & \textbf{Pull-up (GPIO 32)} & \textbf{Pull-down (GPIO 33)} \\ \hline
		Tidak ditekan & HIGH (TIDAK DITEKAN) & LOW (TIDAK DITEKAN) \\ \hline
		Ditekan & LOW (DITEKAN) & HIGH (DITEKAN) \\ \hline
	\end{tabular}
\end{table}
Pada pull-up, resistor menarik pin ke 3,3~V saat tombol terbuka, dan tombol menghubungkan pin ke GND saat ditekan, sehingga logikanya aktif LOW. Pada pull-down, resistor menarik pin ke GND saat tombol terbuka, dan tombol menghubungkan pin ke 3,3~V saat ditekan, sehingga logikanya aktif HIGH. Kode memakai \texttt{stateUp == LOW} dan \texttt{stateDown == HIGH} untuk menyesuaikan perbedaan logika ini.

Pada percobaan awal serial monitor menampilkan GPIO 32 (pull-up) selalu DITEKAN walaupun tombol tidak disentuh, artinya pin terbaca LOW terus. Hal ini konsisten dengan pin yang tidak tertarik ke 3,3~V dan/atau tersambung ke GND akibat kesalahan wiring, sehingga rangkaian tidak berfungsi sebagai pull-up. Asisten praktikum memperbaiki wiring tersebut tanpa penjelasan rinci, dan kami tidak mencatat bagian mana yang diubah, sehingga penyebab pastinya tidak dapat kami tentukan. Setelah wiring diperbaiki, GPIO 32 (pull-up) dan GPIO 33 (pull-down) terbaca TIDAK DITEKAN saat kedua tombol tidak ditekan, sesuai Tabel~\ref{tab:pullup} (tangkapan layar pada Langkah Percobaan 3, Bagian A). Pembacaan saat tombol ditekan pada Bagian A tidak terdokumentasi dengan bukti yang dapat dipastikan, sehingga kondisi DITEKAN yang benar tidak dapat kami tunjukkan dari tangkapan layar.

\foto[0.42]{p3_wiring_salah}{Rangkaian pada trainer board saat GPIO 32 terbaca DITEKAN terus (sebelum wiring diperbaiki)}
\foto[0.9]{p3_serial_salah}{Serial monitor saat GPIO 32 (pull-up) terbaca DITEKAN terus walaupun tombol tidak ditekan}

Pull-up dan pull-down diperlukan agar pin input tidak \emph{floating}. Pin yang mengambang memiliki impedansi sangat tinggi sehingga levelnya tidak terdefinisi dan dipengaruhi noise di sekitarnya. Kejadian di atas menunjukkan bahwa ketika jalur pull-up tidak benar, pembacaan bisa salah secara konsisten, bukan hanya acak.

Resistor eksternal 10~k$\Omega$ dan resistor internal memiliki fungsi yang sama tetapi berbeda dalam praktik. Dengan 3,3~V, arus ke GND saat tombol ditekan pada pull-up 10~k$\Omega$ adalah $3{,}3/10.000 = 0{,}33$~mA. Resistor pull internal ESP32 bernilai sekitar 45~k$\Omega$ (nilai tipikal pada datasheet), sehingga arusnya sekitar 0,07~mA. Resistor internal menghemat komponen dan kabel, tetapi tarikannya lebih lemah sehingga lebih rentan noise pada kabel yang panjang. Resistor eksternal lebih kuat dan nilainya dapat dipilih.
\ifBagianB
Pada Kode Program B (pin dan tombol yang sama), serial monitor menampilkan GPIO 32 (pull-up internal) TIDAK DITEKAN dan GPIO 33 (pull-down internal) TIDAK DITEKAN saat tombol tidak disentuh, lalu GPIO 33 berubah menjadi DITEKAN saat tombol pull-down ditekan. Pola ini sama dengan Kode Program A dan sesuai teori bahwa \texttt{INPUT\_PULLUP} dan \texttt{INPUT\_PULLDOWN} mengaktifkan resistor yang berfungsi sama dengan resistor eksternal, sehingga rangkaian tetap bekerja walaupun resistor 10~k$\Omega$ dilepas. Tangkapan layar yang kami miliki tidak memuat pembacaan GPIO 32 (pull-up) saat ditekan, sehingga bagian itu tidak dapat kami dokumentasikan.
\fi

\subsection{Percobaan 4 dan 5}
Kedua percobaan tidak dilaksanakan sehingga tidak ada data yang dianalisis. Pengaruh nilai waktu debounce (10, 50, dan 200~ms) dan pembacaan ADC dengan LDR belum kami uji langsung. Teori debouncing dan ADC baru diterapkan pada Tugas Modul, itu pun dalam simulasi.

%===========================================================%
\section{Hasil Tugas Modul}
\subsection{Rancangan}
Tugas modul meminta simulasi matriks pushbutton 3x3 yang mengendalikan matriks LED 3x3, dengan tombol pada posisi (baris, kolom) menyalakan LED pada posisi yang sama dan wajib memakai debouncing. Simulasi dibuat di Wokwi dengan ESP32 DevKit. Matriks tombol dibaca dengan teknik \emph{scanning}: baris menjadi output dan kolom menjadi input dengan pull-up internal. Matriks LED dibuat dari sembilan LED yang masing-masing memakai satu pin GPIO dan resistor 220~$\Omega$, dengan susunan posisi yang sama seperti tombol. Pemetaan pin ditunjukkan Tabel~\ref{tab:pinmatrix}.
\begin{table}[H]
	\centering
	\caption{Pemetaan pin matriks 3x3}
	\label{tab:pinmatrix}
	\begin{tabular}{|l|l|l|l|}
		\hline
		\textbf{Fungsi} & \textbf{Baris/Kolom 1} & \textbf{Baris/Kolom 2} & \textbf{Baris/Kolom 3} \\ \hline
		Baris tombol (output) & GPIO 16 & GPIO 17 & GPIO 18 \\ \hline
		Kolom tombol (input pull-up) & GPIO 19 & GPIO 21 & GPIO 22 \\ \hline
		LED baris 1 & GPIO 13 & GPIO 14 & GPIO 25 \\ \hline
		LED baris 2 & GPIO 26 & GPIO 27 & GPIO 32 \\ \hline
		LED baris 3 & GPIO 33 & GPIO 4 & GPIO 23 \\ \hline
	\end{tabular}
\end{table}
Pada setiap putaran scanning, satu baris ditarik ke LOW sementara baris lain dibiarkan \emph{high-Z} (mode input). Jika tombol pada baris aktif ditekan, kolomnya terbaca LOW. Debouncing dilakukan per tombol: pembacaan mentah dibandingkan dengan pembacaan sebelumnya, dan setiap perubahan mengulang penghitungan waktu. Status tombol baru diterima dan LED diubah hanya setelah pembacaan stabil selama \texttt{DEBOUNCE\_MS} = 30~ms. LED mengikuti status tombol yang sudah stabil, yaitu menyala selama tombol ditekan.
\foto{tm_rangkaian_wokwi}{Rangkaian simulasi Wokwi saat tidak ada tombol ditekan: ESP32, matriks pushbutton 3x3, dan matriks LED 3x3 (seluruh LED mati)}

\subsection{Kode Program}
\lstinputlisting[language=C++, caption={Kode program matriks pushbutton 3x3 dengan debouncing}]{Tugas_Modul/main.cpp}

\subsection{Hasil Simulasi}
Simulasi dijalankan di Wokwi dan kesembilan tombol diuji satu per satu dengan cara ditekan lalu dilepas. Pada setiap pengujian, LED pada posisi (baris, kolom) yang sama dengan tombol menyala selama tombol ditekan dan padam ketika tombol dilepas, sedangkan delapan LED lainnya tetap mati. Seluruh sembilan pasangan tombol dan LED sesuai dengan pemetaan pada Tabel~\ref{tab:pinmatrix}. Serial monitor mencatat setiap perubahan status, yaitu satu baris \texttt{DITEKAN} saat tombol ditekan dan satu baris \texttt{DILEPAS} saat dilepas, tanpa baris ganda untuk satu kali penekanan.

Batasan pengujian: tombol hanya ditekan satu per satu. Penekanan beberapa tombol sekaligus tidak diuji. Selain itu, pushbutton pada simulator tidak memodelkan pantulan kontak mekanik, sehingga manfaat debouncing tidak dapat diamati langsung dalam simulasi ini dan hanya didasarkan pada logika program serta penjelasan pada jawaban analisis.
\foto[0.6]{tm_sim_semua}{Matriks LED pada simulasi untuk kesembilan tombol yang ditekan satu per satu: hanya LED dengan posisi yang sama yang menyala}
\foto{tm_sim_tombol}{Simulasi dengan tombol baris 1 kolom 1 ditekan: LED baris 1 kolom 1 menyala dan serial monitor mencatat kejadian}
\foto{tm_sim_serial}{Keluaran serial monitor saat tombol ditekan dan dilepas}

\subsection{Jawaban Pertanyaan Analisis}
\begin{enumerate}
	\item \textbf{Mengapa debouncing tetap diperlukan?} Bouncing adalah sifat fisik setiap kontak mekanik, bukan akibat banyaknya tombol yang ditekan. Pada saat kontak tersambung atau terlepas, pegas logam memantul beberapa kali dalam waktu singkat (orde milidetik) sehingga sinyal bolak-balik antara HIGH dan LOW sebelum stabil, seperti dijelaskan pada dasar teori debouncing. Matriks 3x3 hanya mengatur cara pembacaan, sedangkan setiap tombol tetap berupa satu kontak mekanik yang memantul. Loop scanning membaca tombol sangat cepat, jauh lebih cepat daripada durasi pantulan, sehingga satu kali penekanan akan terbaca sebagai beberapa perubahan. Tanpa debouncing, LED berkedip atau berubah status mengikuti pantulan, dan pada program yang menghitung kejadian satu penekanan dihitung berkali-kali. Dengan menunggu sinyal stabil selama 30~ms, pantulan diabaikan dan LED hanya berubah satu kali per penekanan. Gangguan singkat (\emph{glitch}) yang lebih pendek dari 30~ms juga tidak diterima.

	\item \textbf{Konfigurasi pull-up pada pembacaan kolom.} Kolom dikonfigurasi \texttt{INPUT\_PULLUP}, yaitu resistor internal menarik pin kolom ke 3,3~V. Saat tombol tidak ditekan, pin kolom terbaca HIGH. Saat tombol ditekan dan barisnya sedang di-scan (LOW), kolom tersambung ke baris sehingga terbaca LOW, jadi logikanya aktif LOW. Kondisi HIGH ini tidak boleh dibiarkan \emph{floating} karena pin input CMOS memiliki impedansi sangat tinggi, sehingga levelnya tidak terdefinisi dan mudah berubah akibat noise, induksi dari kabel, atau muatan liar. Pembacaan yang tidak menentu membuat tombol terdeteksi tertekan padahal tidak, sehingga LED berkedip sendiri. Hal serupa kami amati pada Percobaan 3: ketika jalur pull-up tidak terpasang dengan benar, GPIO 32 terbaca LOW terus dan tombol dianggap selalu ditekan.

	\item \textbf{Digital input dan analog input (ADC).} Pada digital input, mikrokontroler hanya membandingkan tegangan pin dengan ambang logika dan menghasilkan dua kondisi, HIGH atau LOW. Pada ESP32 ambang tersebut sekitar 0,25 dan 0,75 dari tegangan catu (nilai kira-kira dari datasheet), dan tegangan di antaranya tidak memberi informasi tambahan. Pada analog input, ADC mengukur besar tegangan dan mengubahnya menjadi bilangan. ADC 12-bit pada ESP32 membagi rentang sekitar 0 sampai 3,3~V menjadi $2^{12} = 4096$ tingkat (0 sampai 4095), dengan resolusi sekitar $3{,}3/4095 \approx 0{,}81$~mV per tingkat. Bila sebuah LED diganti LDR dalam rangkaian pembagi tegangan dengan resistor 10~k$\Omega$, tegangan keluaran berubah kontinu mengikuti hambatan LDR, yang bergantung pada intensitas cahaya. ADC memetakan tegangan itu ke salah satu dari 4096 nilai, sehingga dapat menyatakan banyak tingkat terang, bukan hanya terang atau gelap. Perlu dicatat bahwa ADC ESP32 tidak sepenuhnya linear di dekat ujung rentang.
\end{enumerate}

%===========================================================%
\section{Kesimpulan}
Praktikum Modul 1 berhasil dijalankan sampai Percobaan 3. Program Blink berjalan pada STM32 Blackpill dengan framework Arduino melalui ST-Link (LED PC13 aktif LOW) dan pada ESP32 dengan framework ESP-IDF (LED GPIO 2). Kedua percobaan sesuai teori: periode kedipan 1000~ms mengikuti dua \texttt{delay} 500~ms. Kendala upload pada ST-Link ternyata berasal dari sambungan fisik, bukan dari kode, dan pada ESP-IDF path project harus tanpa spasi serta file program harus berekstensi \texttt{.c}.

Pada Percobaan 3, pembacaan saat tombol tidak ditekan sesuai teori (pull-up HIGH dan pull-down LOW), dan kesalahan wiring terbukti dapat membuat pembacaan salah secara konsisten (GPIO 32 selalu terbaca tertekan). Secara teori pull-up berlogika aktif LOW dan pull-down aktif HIGH.
\ifBagianB
Pada kondisi yang terdokumentasi, resistor internal ESP32 (sekitar 45~k$\Omega$) menghasilkan pembacaan yang sama dengan resistor eksternal 10~k$\Omega$, dengan perbedaan pada kekuatan tarikan dan jumlah komponen.
\else
Resistor eksternal 10~k$\Omega$ dan resistor internal ESP32 (sekitar 45~k$\Omega$) secara teori berfungsi sama, dengan perbedaan pada kekuatan tarikan dan jumlah komponen. Hal ini belum kami uji langsung karena Bagian B tidak dilaksanakan.
\fi

Percobaan 4 dan 5 tidak dilaksanakan karena waktu praktikum habis, sehingga kesimpulan mengenai pengaruh nilai debounce dan pembacaan ADC tidak dapat diambil dari percobaan langsung. Pada Tugas Modul, simulasi di Wokwi menunjukkan bahwa scanning baris-kolom dengan pull-up internal dapat memetakan kesembilan tombol ke kesembilan LED dengan benar, dengan logika debouncing per tombol berjendela stabil 30~ms. Pengaruh debouncing terhadap pantulan kontak belum teramati karena simulator tidak memodelkan pantulan, dan hasil simulasi ini belum diuji pada perangkat keras.

%===========================================================%
\section{Lampiran}
\subsection{Dokumentasi saat praktikum}
\foto{lamp_praktikum_1}{Dokumentasi praktikan saat melaksanakan praktikum Modul 1}
\foto{lamp_praktikum_2}{Dokumentasi rangkaian dan perangkat yang digunakan saat praktikum Modul 1}
% Bila kedua foto Lampiran belum ada (mode final), tampilkan keterangan singkat.
\ifShowMissing\else
\IfFileExists{P1/img/lamp_praktikum_1.png}{}{\IfFileExists{P1/img/lamp_praktikum_1.jpg}{}{%
\IfFileExists{P1/img/lamp_praktikum_2.png}{}{\IfFileExists{P1/img/lamp_praktikum_2.jpg}{}{%
Foto dokumentasi kegiatan praktikum tidak tersedia. Bukti pelaksanaan percobaan berupa tangkapan layar pada bagian Langkah-Langkah Percobaan dan Hasil Tugas Modul.}}}}
\fi
